%----------------------------------------------------------------------
\chapter{Software top-level architectural design}
\label{ch:toplevel}
%----------------------------------------------------------------------

%----------------------------------------------------------------------
\section{Overall architecture}
\label{sc:architecture}
%----------------------------------------------------------------------

\ssat{} provides a framework for computing and analyzing stack usage
for the different subprograms in the partition and extracting worst case
prior to execution time. It is made up by the following basic components:

\begin{itemize}

\item Compiler

\item Call graph generator and decorator

\item Worst case analyzer

\item GNAT Programming Studio (GPS)

\end{itemize}

This chapter will focus on describing the capabilities required for these
components.

%----------------------------------------------------------------------
\section{Data representation}
\label{sc:data}
%----------------------------------------------------------------------

Two things are required for analyzing statically stack usage: the call
graph information, and the stack requirements for each
subprogram. Once we have that information, we can traverse the
decorated call graph to look for the worst case path according to its
stack requirements.

First we need to define how to represent the call graph corresponding
to the application. We do that in the form of a directed call graph,
in which nodes will represent subprograms that are defined or called,
and directed edges between nodes will represents potential calls. The
graph needs to be able to store the stack requirements for each
subprogram, in terms of boundaries and attributes for these
boundaries.

There are four types of nodes (see figure \ref{fg:node-hierarchy}),
representing the different types of subprograms that may be called:

\begin{itemize}

\item Explicit node. It corresponds to a subprogram that has been
  found and analyzed in the set of sources under analysis, and can be
  unambiguously identified at compile-time.

\item Dispatching node. This is a node representing a dispatching
  call so we do not know statically the actual subprogram that will
  be called. It uses the dynamic dispatching capability provided by
  class-wide type. We use this node that will then point to every
  known operation to which it can dispatch.

\item Indirect node. It represents a node that is a dereference of
  access-to-subprogram so we do not know statically the concrete
  subprogram that will be called.

\item Unresolved node. It corresponds to a subprogram that is
  referenced from the set of sources being analyzed, but it is not
  part of these sources (it may be a call to an external library which
  does not provide the required call graph and stack usage
  information). Hence, there is no information about its stack usage
  nor about subprograms that are called from this node.

\end{itemize}

\begin{figure}
\centerline{\includegraphics*[scale=0.6]{node-hierarchy.eps}}
\caption{Node class hierarchy}
\label{fg:node-hierarchy}
\end{figure}

Dispatching and indirect are considered as \emph{fake} nodes because
they do not represent real subprograms. They represent calls whose
target subprogram can only be determined at run time. Representing
these entities as intermediate fake nodes (pointing to all possible
target subprograms) is practical for representing and analyzing the
call graph.

As it was said previously, these nodes need to store the stack
requirements for the subprogram they represent, and this information
depends on the node they represent. The information that we need to
store about explicit nodes is:

\begin{itemize}

\item Subprogram name. This is a \prog{String} that contains the
  symbol name for the node.
\item Location list. It is a list of triplets file name, line number,
 and column number, describing the place where the subprogram body
 appears. The list of locations will typically contain just one entry,
 and it will contain a list of more than one location in the case of
 instantiations of generic units.

\item Stack usage information. It contains stack requirements for the
  represented subprogram, in terms of boundaries and attributes for
  these boundaries. Attributes indicate whether the stack usage for
  the subprogram can be statically known at compile-time, or partially
  known (there is a dynamic part that depends on data arguments passed
  to the frame, which can be bounded or unbounded).

\item Accumulated stack usage information. It contains stack
  requirements for the complete call tree starting in the referred
  subprogram, with the corresponding attribute.

\item Target nodes. It is the set of possible nodes that can be called
  from this subprogram, and it is represented as a list of directed
  edges (see definition bellow) between subprograms. These calls are
  ordered according to their accumulated stack requirements.

\item Analysis information. This is a \prog{Boolean} variable that
  indicates whether the subprogram is part of the path being analyzed
  for obtaining the worst case stack usage. It is very practical for
  detecting and handling cycles in the call graph.


\item Cycle information. A subprogram may be part of different cycles
  in the call graph, and this field contains the set of cycles to
  which the referred subprogram belongs.

\end{itemize}

All the rest of nodes contain the same information except that about
local stack usage because it is meaningless in other contexts.

Directed edges are used for representing \emph{may call}
relationships between subprograms. For each of these edges we store
the following information:
 
\begin{itemize}

\item Target node. Pointer to the node that is being called.

\item Location list. It is a list of triplets file name, line number,
 and column number, describing the place where the call
 appears. The list of locations will typically contain just one entry,
 and it will contain a list of more than one location in the case of
 instantiations of generic units.

\item Feasibility. This flag is placed here for convenience to be able
  to incorporate information provide by the user about calls that are
  known not to happen (because of being dead or deactivated code, for
  example).

\item Analysis information. This is a \prog{Boolean} variable that
  indicates whether the call is part of the path being analyzed
  for obtaining the worst case stack usage. It is very practical for
  detecting and handling cycles in the call graph.

\end{itemize}

The basic data types and subprograms required for static stack
analysis is represented in the following package spec:

\lgrindfile{static_stack_analysis}

%----------------------------------------------------------------------
\section{Compiler}
%----------------------------------------------------------------------

\ssat{} is based on the development of  dedicated compiler extensions
to produce specialized outputs for stack requirements analysis
purposes.

One limitation is that a compiler cannot provide information on elements it
does not process, such as COTS operating system services for which sources are
not available or very low level routines developed in assembly language.

When worst case bounds are a strong concern, not having the sources of some
components is rare, however, and the stack usage in assembly routines is
usually simple enough to be accounted for manually.

The compilation process may also not be able to grasp the interrupt handling
bits necessary to size the worst case amount of interrupt related stack, be it
for hardware interrupt or signal handlers.

Interrupt handling always requires very careful design and coding, though, so
the information could at least also be provided to the framework by the user,
or accounted for separately.

Leveraging on a compiler's internals has a number of advantages:

\begin {itemize}
\item {\it Reduced effort for new target architectures}:

A compiler typically knows everything about stack allocations for the code it
processes.

When already structured to support multiple target architectures, this
knowledge may be used for each and alleviates the need of an associated
comprehensive machine instruction parser, potentially difficult and
error-prone for complex processors.

\item {\it User level guards}:

A compiler can be tailored with options to raise warnings or errors on user
level constructs known to trigger stack allocation patterns that may cause
troubles to a stack analysis framework involved later on.

This points directly at the user level construct, and so very early in the
development process, both being of precious value.

\item {\it Access to high level information}:

A compiler has visibility on semantic information that can help tackle some of
the challenging issues we have previously identified.

Consider an Ada \verb|Integer| subtype with range \verb|1..5| for
example.

If a variable of this subtype is used to size a local array, the range
knowledge may be used to compute a bound on the array size, and so on the
corresponding stack allocation.

Potential targets of an indirect calls are another example.

Based on subprogram profiles and actual references to subprograms, the
compiler can provide a limited but exhaustive list of subprograms possibly
reached by an indirect call.

In both cases, the compiler information at hand is extremely hard, if not
impossible, to infer from a machine level representation.

\item {\it Scalability}:

Support for stack usage outputs is unlikely to change how a compiler scales up
against application sizes, so a compiler-based stack analysis component will
scale up as well as the initial compiler did.

Besides, a compiler has the opportunity to output no more than what is really
relevant for later analysis purposes, which makes this output easier to digest
than a huge machine level representation of the code.

\end{itemize}

%----------------------------------------------------------------------
\subsection{Data generated by the compiler}
%----------------------------------------------------------------------

The compiler can generate a file per compilation unit (controlled by a
command line option, \prog{-fcallgraph-info=su}), that contain the
required information for generating the call graph, and decorating it
with stack usage information.

The extension \prog{.ci} (call graph information) will be used for the
compiler-generated files containing this information (one file per
compiled unit).

%----------------------------------------------------------------------
\subsubsection{Call graph information}
\label{compiler-call-graph}
%----------------------------------------------------------------------

The call graph is made up by nodes and directed edges between
nodes. The compiler produces this information on a per-unit basis in
the form of a text file in VCG \refcite{VCG} graph description
language.

The types of nodes that are generated by the compiler are:

\begin{itemize}

\item \emph{Defined-in} node. It corresponds to a subprogram defined
  in the unit being compiled.

\item \emph{Referenced-from} node. It identifies a subprogram
  referenced from the unit being compiled but it is either defined in a
  different compilation unit or not known at compile time. These nodes
  can be categorized into:

  \begin{itemize}
  
  \item \emph{Proxy} node. It represents a subprogram being called
  from this unit and known to be defined in a different compilation
  unit.

  \item \emph{Dummy proxy} node. It corresponds to a subprogram which
  is called from the unit being compiled but its target not statically
  known. It corresponds to either an indirect call or a dispatching
  call.

  \end{itemize}

\end{itemize}

All the nodes, except \emph{dummy proxy} nodes, are represented by:

\begin{itemize}

\item Title. It contains the full name of the subprogram (symbol
  name).

\item Label. Ada/C name of the subprogram plus a list of locations (file
 name and line number). The list of locations will typically contain
 just one entry, and it will contain a list of more than one location
 in the case of instantiations of generic units.

\end{itemize}

\emph{Dummy proxy} nodes simply contain a mark indicating that they
are not known at compile time.

Note that the nodes that appear are those for which there is an
explicit call. Inlined subprograms, for example, will not appear in
the call graph since no actual call is made.

For every compiled unit, a list of directed edges will be generated
corresponding to every call appearing in the unit being compiled. They
materialize a {\em may\_call} relationship, where the source
subprogram {\em may\_call} the destination subprogram.These edges are
represented by:

\begin{itemize}

\item Source. This is the title corresponding to the calling
  subprogram. It can only be a \emph{defined-in} node.

\item Target. Title corresponding to the called subprogram. It can be
  any type of node.

\end{itemize}

%----------------------------------------------------------------------
\subsubsection{Stack usage information}
\label{compiler-stack-usage}
%----------------------------------------------------------------------

The compiler generates stack usage information for each subprogram
defined in the compilation unit. The following information will be
generated for each subprogram, and will later be used for decorating
the call graph tree:

\begin{itemize}

\item List of locations (file name and line number) where the call
 appear. The list of locations will typically contain  just one entry,
 and it will contain a list of more than one location in the case of
 instantiations of generic units.

\item Ada/C name of the subprogram.

\item Stack requirements. Maximum computable amount of stack the
  subprogram may ever allocate on the stack. It corresponds only to
  the amount that  can be statically bounded.

\item Attributes for the stack requirements. It indicates the stack
  allocation pattern, which may be:

  \begin{itemize}

  \item Static. Only constant allocations are made, so the defined
    stack requirement.

  \item Dynamic and bounded. Dynamic allocations may occur in addition
  to constant ones, but the sum of all these allocations is never
  greater than the defined stack requirement.

  \item Dynamic and unbounded. Dynamic allocations may occur in
  addition to constant ones, and the sum of all these allocations can
  be greater than the defined stack requirement to an unknown extent.

  \end{itemize}
  
\end{itemize}

%----------------------------------------------------------------------
\section{Call graph generator and decorator}
%----------------------------------------------------------------------

This component is in charge of generating the complete call graph for
every entry point in the partition using the partial call graph
information previously generated by the compiler. The stack usage
information corresponding to the subprogram, including the information
about whether these values are static, dynamic, or bounded, is also
part of the call graph information.

Some entry points are automatically identified (main entry point, task
bodies, interrupt handlers). The user can also specify any additional
entry point that needs to be considered in the analysis. These entry
points will be stored in a dynamic list.

The compiler cannot always determine the required information for a
complete analysis. Potential sources for incomplete information to be
addressed are:

\begin{itemize}

\item Indirect calls. These calls correspond to dereference of
  access-to-subprogram, so we do not know statically the actual
  subprogram that will be called because it depends on the data flow.

\item Dispatching calls. These calls use the dynamic dispatching
  capability provided by class-wide type, so we do not know statically
  the actual subprogram that will be called because it depends on the
  actual object used.

\item Dynamic unbounded frames. There is a dynamic part in the stack
  requirements that depends on data arguments passed to the frame,
  which cannot be statically bounded.

\item External calls. There is no stack usage information about the
  target call.

\item Cycles. In presence of control flow cycles involving stack
  consumption, the worst case bound is a function of the maximum
  number of cycle iterations.

\end{itemize}

The \emph{graph generator and decorator} can help the user (or an
external tool) to fill the required information; it will generate a
file containing the pattern to fill the missing pieces of
information. The user will then be able to provide more precise
information about potential target addresses and stack usage in what
is called \emph{user-defined call graph and stack usage files}
(\prog{.ciu}).

The format of the compiler generated file has been defined in section
\ref{compiler-call-graph}. The format of the \emph{user-defined
  call graph and stack usage files} will be a simplified version of
the data generated by the compiler in order to simplify the design of
the of parser.

The \emph{call graph generator and decorator} will parse both
compiler-generated and user defined files and will create a call graph
tree made up by a list of nodes and edges as described in \ref{sc:data}.

%----------------------------------------------------------------------
\section{Worst case analyzer}
%----------------------------------------------------------------------

Once we have the annotated call tree we need to look for the longest
paths in terms of stack usage. We will start on every entry point and
fill the accumulated stack usage, together with the pointer to the
longest path, in a recursive way. The pseudocode for this recursive
function is the following:
\begin{alltt}
   function Get_Accumulated_Stack_Usage (Node) return Stack_Usage is
     Tmp_Accumulated  := null;
     New_Accumulated  := null;
     Tmp_Longest_Path := null;
   begin

     if not Cycle_Detected and then Node.Accumulated_Stack_Usage = null then
       --  Compute the information only if it is not yet available

       for Target_Node in List_Of_Target_Nodes loop
         New_Accumulated :=
           Get_Accumulated_Stack_Usage (Target_Node);

         if New_Accumulated > Tmp_Accumulated then
           Tmp_Accumulated  := New_Accumulated;
           Tmp_Longest_Path := Target_Node;
         end if;
       end loop;

       Node.Accumulated_Stack_Usage := Tmp_Accumulated;
       Node.Longest_Path := Tmp_Longest_Path;
     end if;

     return Node.Accumulated_Stack_Usage;

   end Get_Accumulated_Stack_Usage;


   for Node in List_Of_Entry_Points loop
     Node.Accumulated_Stack_Usage := Get_Accumulated_Stack_Usage (Node);
   end loop; 
\end{alltt}

% Note that the existence of a node with a unknown or dynamic node make
% it the longest path.

%----------------------------------------------------------------------
\section{GNAT Programming Studio (GPS)}
%----------------------------------------------------------------------

TBD
